While GraphRAG-Router demonstrates strong performance on six QA benchmarks, we acknowledge a few limitations that present opportunities for future work. First, the current study has a limited task scope and primarily considers Wikipedia-based QA settings. it remains an open question whether the proposed GraphRAG-Router can generalize equally well to more complex retrieval-heavy tasks, such as multi-hop reasoning over long enterprise/scientific documents. Second, our current evaluation is conducted mainly in an offline retrieval setting, where the underlying corpora, graph indices, and candidate GraphRAG systems are pre-built and fixed. It remains unclear whether GraphRAG-Router can maintain the same effectiveness in more dynamic online environments with continuously updated knowledge sources and evolving retrieval indices.